% concatenated from main.tex, abstract.tex, introduction.tex, preliminary.tex, mainresult.tex, proof.tex, numerical.tex, appendix.tex

\documentclass[a4paper,fleqn]{cas-sc}
\usepackage[authoryear]{natbib}

\usepackage{amsthm,amsmath,amssymb}
\usepackage{mathtools,setspace,enumitem}

\theoremstyle{plain}
\newtheorem{theorem}{Theorem}[section]
\newtheorem{remark}[theorem]{Remark}
\newtheorem{claim}[theorem]{Claim}
\newtheorem*{theorem*}{Theorem}
\newtheorem*{claim*}{Claim}
\newtheorem*{remark*}{Remark}
\newcommand{\R}{\mathbb{R}}
\newcommand{\E}{\mathbb{E}}
\newcommand{\tX}{\tilde{X}}
\newcommand{\tY}{\tilde{Y}}
\newcommand{\tE}{\tilde{E}}
\newcommand{\p}{\mathcal{P}}
\newcommand{\g}{\mathcal{G}}
\newcommand{\tSigma}{\tilde{\Sigma}}
\newcommand{\tL}{\tilde{L}}
\newcommand{\tb}{\tilde{\beta}}
\newcommand{\e}{\varepsilon}
\newcommand{\la}{\lambda}
\newcommand{\be}{\beta}
\newcommand{\ga}{\gamma}
\newcommand{\te}{\theta}
\newcommand{\hte}{\hat{\theta}}
\newcommand{\tte}{\tilde{\theta}}
\newcommand{\mtr}{\mathrm{tr}}


\begin{document}
\let\WriteBookmarks\relax
\def\floatpagepagefraction{1}
\def\textpagefraction{.001}

\title[mode=title]{Bidirectional Random Projections} 
\shorttitle{Bidirectional Random Projections}
\shortauthors{}

\author[1]{Chao Lan}
\ead{clan@ou.edu}
\author[1]{Luyuan Yang}
\affiliation[1]{
organization={
School of Computer Science, 
University of Oklahoma}, 
country = {U.S.}
}

\begin{abstract}
This paper analyzes  
bidirectional random 
projections for ordinary 
least squares (OLS) regression 
under the fixed design setting. 
Let $(X,Y) \in \R^{n \times p} 
\times \R^n$ be a sample and 
$R \in \R^{n_1 \times n}, 
W \in \R^{p \times p_1}$ be 
two properly distributed 
random projections. 
We develop an expected excess 
loss bound for the OLS estimator 
built on $(WXR, WY)$. Compared 
to an established bound for 
OLS estimator built on $(XR, Y)$, 
the gap is approximately $O\left( p_1 
+ C \frac{1}{p_1} \right)$, 
where $C$ scales with $n_1/n$ 
and can be negative for small $n_1/n$. 
Its implications are confirmed 
by numerical results on 
real-world data. 
\end{abstract}

\begin{keywords}
Random Projection \sep 
Bidirectional Random Projection \sep 
Ordinary Least Square \sep 
Expected Excess Loss \sep 
High Dimensional Data 
\end{keywords}


\maketitle


\section{Introduction}

It is well-known that 
ordinary least squares (OLS) regression 
can be scaled up for high-dimensional data 
by using random projection to reduce 
data dimension. Excess loss for 
the projected OLS (pOLS) estimator has been well studied 
\citep{maillard2009compressed,fard2012compressed,kaban2014new,slawski2017compressed}. 
On the other hand, random projection 
can also be used to reduce data number 
for faster estimation 
\citep{zhou2007compressed}. As a result, 
when working with large-scale high 
dimensional data, it may be desirable to 
apply random projection in both directions. 
However, statistical properties of such 
bi-directionally projected OLS (bOLS) 
estimator remains unknown. 

This paper develops an excess 
loss bound for bOLS estimator. 
By inspecting its gap to the 
bound for pOLS estimator, 
we obtain three implications: 
(i) bOLS performs worse than pOLS; 
(ii) as the projected data dimension increases, their gap 
first decreases and then increases, and the pivot happens earlier 
when the reduced data number is smaller; 
(iii) as the projected data dimension 
increases, both losses first decrease 
and then increase, but bOLS's 
loss bounces back earlier. These 
theoretical implications are confirmed 
by numerical results on real-world data. 

\section{Preliminaries}

Consider the fixed design setting. 
Let $x_1, \ldots, x_n \in \R^p$ be 
$n$ fixed points and $y_1, \ldots, 
y_n$ be their labels generated by 
$y_i = x_i^T \be_* + \e_i$, 
where $\be_* \in R^p$ is fixed 
and $\e_i$'s are independent 
random noises with $\E \e_i = 0$ 
and $var(\e_i) = \sigma^2$ for 
some fixed $\sigma >0$.\ \  
Let $X \in \R^{n \times p}$ 
be a sample matrix with 
$x_i^T$ on the $i_{th}$ row, 
$Y \in \R^n$ be a label 
vector with $Y_i = y_i$ 
and $E \in \R^n$ be a 
noise vector with $E_i = \e_i$. 
Then, the above setting 
can be expressed as 
\begin{equation}
Y = X \be_* + E.    
\end{equation}

Let $\p(v^2)$ denote a Gaussian 
distribution on $\R$ with zero mean 
and variance $v^2$. 
Consider two random projections
$W \in \R^{n_1 \times n}$ 
and $R \in \R^{p \times p_1}$. 
For $W$, assume it has i.i.d. 
entries following 
\begin{equation}
\label{eq:sampleW}
W_{ij} \sim \p(1/n).      
\end{equation}
For $R$, assume it has i.i.d. 
entries following 
\begin{equation}
\label{eq:sampleR}
R_{ij} \sim \p(1/p_1).     
\end{equation}
Consider two OLS estimators. 
Let $\hat{\te}_R \in R^{p_1}$ 
be a projected estimator built 
on $(XR, Y)$ i.e., 
\begin{equation}
\label{eq:def_teR}
\hat{\te}_R = \arg\min_{\te} 
\frac{1}{n} \parallel X R \te - Y \parallel^2,      
\end{equation}
where $||\cdot||$ is $L_2$ norm. 
Let $\hat{\te}_{RW}$ be a bi-directionally 
projected estimator built on $(WXR, WY)$ i.e., 
\begin{equation}
\hat{\te}_{RW} = \arg\min_{\te} 
\frac{1}{n_1} \parallel W X R \te - W Y \parallel^2.
\end{equation}

We are interested in comparing  
$\hat{\te}_{RW}$ with $\hat{\te}_R$. 
Comparison will be based on 
two loss functions: the 
expected square loss of 
any estimator $\be \in \R^p$
is defined as 
\begin{equation}
L(\be) = \E_{Y}\ \frac{1}{n} 
\parallel X \be - Y \parallel^2. 
\end{equation}
Given $R$, the expected square loss 
of any estimator $\te \in \R^{p_1}$ 
is defined as 
\begin{equation}
L_R(\te) = \E_{Y \mid R}\ \frac{1}{n} 
\parallel X R \te - Y \parallel^2.     
\end{equation}

Finally, 
let $I_s$ be an identity matrix of 
size $s$; for matrix $G$, let 
$||\cdot||_G$ be the induced 
Mahalanobis norm, $\mtr(G)$ be its 
trace, and $\lambda_{g \uparrow}$ 
and $\lambda_{g \downarrow}$ be its 
largest and smallest singular values 
respectively. (Lower case 
of $G$ is used for its singular values.) 
$\E_{Z}$ means an expectation 
is w.r.t. the randomness 
in $Z$, and $\E_{Z \mid Q}$ means 
the expectation is conditioned on $Q$. 
In consecutive arguments, an 
$\E$ without subscript has the 
same subscript as its preceding $\E$.   

\section{Main Result}
\label{sec:mainresult} 

The following is an established
loss bound for $\hat{\te}_R$. 

\begin{theorem}[\cite{kaban2014new}]
\label{thm:rpbound}
Under the fixed design setting, 
for $p_1 < \min(p, n)$: 
\begin{equation}
\E_{R, \hte_R} [L_{R}(\hat{\te}_R)] 
- L(\be_*) \leq \sigma^2 \frac{p_1}{n} 
+ \frac{1}{p_1} || \be_* ||^2_{M}\ 
:= B_R,   
\end{equation}
where $M = \Sigma + (1+\kappa)\, 
\mtr(\Sigma) I_p$ and 
$\Sigma = \frac{1}{n} X^T X$. 
\end{theorem}

Our bound for $\hat{\te}_{RW}$ is 
stated as follows. 

\begin{theorem} 
\label{cor:brpbound}
Under the fixed design setting, 
there exists a 
$W_0 \in \R^{n_1 \times n}$ 
such that for $p_1 < \min(p, n_1)$: 
\begin{equation}
\label{eq:brpbound}
\E_{R, W, \hte_{RW}} 
[L_{R}(\hat{\te}_{RW})] - L(\be_*) 
\leq C_1 \sigma^2 \frac{p_1}{n} 
+ C_2 \frac{1}{p_1} || \be_* ||_{M}^2 
 + b\ := B_{RW},  
\end{equation}
where $M$ is specified in 
Theorem \ref{thm:rpbound}, 
$C_1 = \E_W[\lambda^2_{w \uparrow}]
\lambda^{-2}_{w_0 \downarrow}$,  
$C_2 = \frac{n_1}{n} 
\lambda^{-2}_{w_0 \downarrow}$ 
and $b = (C_2 - 1)L(\be_*)$.
\end{theorem}

It is not hard to see Theorem \ref{cor:brpbound} recovers 
Theorem \ref{thm:rpbound} 
if $W = I_n$ (so that $n_1 = n$ 
and all $\lambda$'s 
equal to 1).
In general, we are interested 
in the gap between two bounds 
and how it behaves as $p_1$ increases. 
To this end, more concrete ideas are needed 
on $\lambda_{w_0 \downarrow}$. 
Since $W$ has i.i.d. entries from 
$N(0,1/n)$, \citep{bai1993limit} showed 
that asymptotically $\lambda_{w \downarrow}$ 
concentrates at 
$1-\sqrt{n_1/n}$ while 
$\lambda_{w \uparrow}$ 
is concentrated at 
$1+ \sqrt{n_1/n}$. 
Then, by the choice of 
$W_0$ in (\ref{eq:road_09}), 
\begin{equation}
\lambda^2_{w_0 \downarrow} 
= \frac{\E_{R, W, \hte_{RW}} 
\lambda^2_{w \downarrow} 
L_{R}(\hte_{RW}) 
}{\E_{R, \hte_{RW}} L_{R}(\hte_{RW})} 
\approx (1 - \sqrt{n_1/n})^2. 
\end{equation}
Let $r = \sqrt{n_1/n}$. The above implies
\begin{equation}
\label{eq:c1}
C_1 = \E_W[\lambda^2_{w \uparrow}] 
\lambda^{-2}_{w_0 \downarrow}  
\approx (\sqrt{n/n_1}+1)^2 
(\sqrt{n/n_1}-1)^{-2} 
= (1+r)^2 (1-r)^{-2}, 
\end{equation}
and 
\begin{equation}
C_2 = \frac{n_1}{n} 
\lambda^{-2}_{w_0 \downarrow}
\approx {(\sqrt{n/n_1}-1)^{-2}} 
= r^2 (1-r)^{-2}.
\end{equation}
Now we can evaluate the gap 
between two bounds. 
It is not hard to see 
\begin{equation}
\label{eq:gap2}
B_{RW} - B_R = (C_1 -1) 
\frac{\sigma^2}{n} p_1 \ + 
\ (C_2-1) || \be_* ||_{M}^2 
\frac{1}{p_1} + b
\approx \frac{4r}{(1-r)^2} 
\frac{\sigma^2}{n}  p_1 
+ \frac{2r - 1}{(1-r)^2} 
||\be_*||_M^2 \frac{1}{p_1} + b.
\end{equation}
The gap is generally positive 
(e.g., for large $r$ or $p_1$),  
which suggests $\hte_{RW}$ performs 
worse than $\hte_{R}$. 

To inspect how the gap behaves 
as $p_1$ increases, we consider 
three scenarios: (i) $p_1$ is 
large, (ii) $p_1$ is small and 
$r$ is large, (iii) $p_1$ is 
small and $r$ is small. 
In (i), the gap is dominated by 
$\frac{2r}{(1-r)^2} \frac{\sigma^2}{n} p_1$ and thus increases 
as $p_1$ grows. In (ii), 
the gap is dominated by 
$\frac{2r - 1}{(1-r)^2} 
||\be_*||_M^2 \frac{1}{p_1}$, 
which has $2r - 1 > 0$ for large 
enough $r$ and thus decreases as $p_1$ grows. In (iii), the gap is 
dominated by the same term, which has $2r - 1 < 0$ for small enough $r$ and thus increases as $p_1$ grows. 
It is noted that, although $r = 0.5$ is the pivot for $\frac{2r - 1}{(1-r)^2} ||\be_*||_M^2 \frac{1}{p_1}$ to transit 
between increasing and decreasing, 
it is not necessarily the pivot for the gap since all constants are 
approximated. Moreover, the pivot for the gap is not necessarily attainable. 
For example, (\ref{eq:gap2}) suggests 
the pivot is approximately at $p_1 = \sqrt{\frac{2r-1}{4r}} 
\frac{\sqrt{n}}{\sigma} ||\be_*||_M$, 
which is only achievable 
if the right side 
is bigger than 1 and $r > 1/2$. Thus, when $r$ is very small (depending on $\sqrt{n} 
||\be_*||_M / \sigma$), 
the gap will only increase 
as $p_1$ increases. 
In summary, as $p_1$ increases, 
the gap generally decrease first 
and then increase; when $r$ is very small, the gap may only increase 
as $p_1$ increases. 

Finally, we inspect the 
pivot of each bound. 
As $p_1$ increases, 
both bounds first decrease 
and then increase. The pivot  
point of $B_R$ is $\sqrt{n ||\be_*||_M^2/\sigma^2}$  
and for $B_{RW}$ it is $C_2/C_1 
\sqrt{n ||\be_*||_M^2/\sigma^2}$. 
It is not hard to see  
$C_2 / C_1 \approx \frac{r^2}{(1+r)^2} 
\leq 1$, which implies the 
loss of $\hte_{RW}$ will bounce back earlier than the loss of $\hte_{R}$. 
\section{Proof of Theorem \ref{cor:brpbound}}

Given $(R,W)$, let the expected 
square loss of $\te$ be  
\begin{equation}
L_{RW}(\te) = \E_{Y \mid R, W}\ 
\frac{1}{n_1} \parallel W X R \te 
- W Y \parallel^2.     
\end{equation} 
Let $\hte_*$ be 
a minimizer of $L_{RW}(\te)$ i.e., 
\begin{equation}
\hte_* = \arg\min_{\te 
\in \R^{p_1}} L_{RW}(\te).
\end{equation}

\subsection{Roadmap}
\vspace{5pt}

We first bound $\E L_{RW}(\hte_{RW})$ 
and then connect it to $\E [L_{R}(\hte_{RW})]$. Given $(R,W)$, it can be 
shown (Claim \ref{claim:02}) 
\begin{equation}
\E_{\hte_{RW} \mid R,W} 
[\hte_{RW}] = \hte_*, 
\end{equation}
which implies (Claim \ref{claim:03}) 
\begin{equation}
\label{eq:road_01}
\E_{\hte_{RW} \mid R,W} [L_{RW}(\hte_{RW})] \leq \E \frac{1}{n_1}
|| W XR \hte_{RW} - W XR \hte_* ||^2 
+ L_{RW}(\hte_*).
\end{equation}
On the right side, the first term 
can be bounded as (Claim \ref{claim:04}) 
\begin{equation}
\label{eq:road_02}
\E_{\hte_{RW} \mid R,W} 
\frac{1}{n_1} 
|| WXR \hte_{RW} - WXR \hte_* ||^2 
\leq \sigma^2 \frac{p_1}{n_1} 
\lambda_{w \uparrow}^2. 
\end{equation}
The second term can be bounded 
as $L_{RW}(\hte_*) \leq 
L_{RW}(R^T\be_*)$ since 
$\hte_*$ minimizes $L_{RW}(\te)$. 
Plugging both back to 
(\ref{eq:road_01}) and 
taking expectation w.r.t. 
$W$ on both sides, we have 
\begin{equation}
\label{eq:road_04b}
\E_{W, \hte_{RW} \mid R} 
[L_{RW}(\hte_{RW})] 
\leq \sigma^2 \frac{p_1}{n_1} 
\E_W [\lambda_{w \uparrow}^2] 
+ \E_W [L_{RW}(R^T\be_*)] 
= \sigma^2 \frac{p_1}{n_1} 
\E [\lambda_{w \uparrow}^2] 
+ L_{R}(R^T\be_*),  
\end{equation}
where the second equality 
is by the fact that 
$\E_W ||W z||^2 
= \frac{n_1}{n} ||z||^2$ for 
any $z$ not depending on $W$. 

For the last term $L_R(R^T \be_*)$, 
its expectation has been bounded 
in \citep{kaban2014new} as 
(Remark \ref{rem:priorkaban})
\begin{equation}
\E_{R} [L_{R}(R^T\be_*)] 
\leq L(\be_*) + \frac{1}{p_1} 
||\be_*||_M, 
\end{equation}
where $M = \Sigma + (1+\kappa)\, 
tr(\Sigma) I_p$ and 
$\Sigma = \frac{1}{n} X^T X$. 
Plugging this in (\ref{eq:road_04b})
after taking expectation 
w.r.t. $R$ on both sides, 
\begin{equation}
\label{eq:road_08}
\E_{R, W, \hte_{RW}} 
[L_{RW}(\hte_{RW})] 
\leq \sigma^2 \frac{p_1}{n_1} 
\E_W [\lambda_{w \uparrow}^2] 
+ \frac{1}{p_1} ||\be_*||_M 
+ L(\be_*). 
\end{equation}
Finally, it can be shown (Claim \ref{claim:05}) 
there exists a $W_0 \in \R^{n_1 \times 
n}$ depending on $R,W,\hte_{RW}$ 
distributions such that 
\begin{equation}
\label{eq:road_09}
\E_{R, W, \hte_{RW}} 
[L_{RW}(\hte_{RW})] 
\geq \E_{R, W, \hte_{RW}} 
[\frac{n}{n_1} 
\lambda^2_{w \downarrow} 
L_{R}(\hte_{RW}) ] = \frac{n}{n_1} 
\lambda^2_{w_0 \downarrow} 
\E_{R, \hte_{RW}} [L_{R}(\hte_{RW})] 
\end{equation}
Combining (\ref{eq:road_08}) 
and (\ref{eq:road_09}) proves 
the theorem. 

\subsection{Supporting Claims}
\vspace{5pt}

We first remark two basic 
observations without proof. 
For any $z \in \R^{n \times k}$ 
(of any $k > 0$), there is 
\begin{equation}
\lambda^2_{w \downarrow} 
||z||^2 \leq ||W z||^2 \leq 
\lambda^2_{w \uparrow} 
||z||^2.    
\end{equation}
In addition, 
the random noise vector 
satisfies $\E_{Y \mid R, W} [E] 
= \vec{0}$ and $\E_{Y \mid R, W} 
[EE^T] = \sigma^2 I_{n}$. 




\begin{claim}
\label{claim:01}
$L_R(\hte_{RW}) \leq \frac{n_1}{n} 
\frac{1}{\lambda_{w \downarrow}^{2}} 
L_{RW}(\hte_{RW})$. 
\end{claim}
\begin{proof}
Since $||W z|| \geq 
\lambda_{w \downarrow} 
||z||$ for any $z$, we have 
\begin{equation}
L_R(\hte_{RW}) = \frac{1}{n}
\E_{Y \mid R, W}||WXR \hte_{RW} - W Y||^2 
\geq \frac{1}{n} \lambda_{w \downarrow}^2  
\E ||XR \hte_{RW} - Y ||^2 
= \frac{n_1}{n} \lambda_{w \downarrow}^2  
\E \frac{1}{n_1} 
||XR \hte_{RW} - Y ||^2. 
\end{equation}
The last expectation is 
$L_R(\hte_{RW})$. This proves the claim.
\end{proof}

\begin{claim}
\label{claim:02}
$\E_{Y \mid R,W} 
[\hte_{RW}] = \hte_*$. 
\end{claim}
\begin{proof}
Let $\tX = WXR$ and $\tY = WY$. 
Since $\hte_{RW}$ minimizes 
$\frac{1}{n_1}||\tX \te - \tY||^2$, 
it is not hard to show 
\begin{equation}
\label{eq:claim02_03}
\hte_{RW} = (\tX^T \tX)^{-1} \tX^T \tY,   
\end{equation}
and thus 
\begin{equation}
\label{claim2:01}
\E_{Y \mid R, W} [\hte_{RW}] 
= (\tX^T \tX)^{-1} \tX^T \E[\tY].
\end{equation}
For $\hte_*$, recall it minimizes 
$L_{RW}(\te)$ and that 
\begin{equation}
L_{RW}(\te) 
= \frac{1}{n_1} 
\E_{Y \mid R, W} 
||\tX \te - \tY||^2
= \frac{1}{n_1} 
\left( ||\tX \te||^2 + 
\E ||\tY||^2 - 2 \E[\tY]^T(\tX \te)\right).
\end{equation}
By taking derivative of the right 
side w.r.t. $\te$ and setting it 
to zero, it is not hard to see 
\begin{equation}
\label{claim2:02}
\hte_* = (\tX^T \tX)^{-1} 
\tX^T \E_{Y \mid R, W} [\tY].  
\end{equation}
Combining (\ref{claim2:01}) 
and (\ref{claim2:02}) proves the claim. 
\end{proof}

\begin{claim}
\label{claim:03}
$\E_{\hte_{RW} \mid R,W} 
[L_{RW}(\hte_{RW})] \leq 
\E \frac{1}{n_1} || W XR 
\hte_{RW} - W XR \hte_* ||^2 
+ L_{RW}(\hte_*)$.
\end{claim}
\begin{proof}
Let $\tX = WXR$ and $\tY = WY$. 
Observe that 
\begin{equation}
\label{eq:claim03:prf_01}
\begin{split}
n_1 \cdot \E_{\hte_{RW} \mid R,W} 
[L_{RW}(\hte_{RW})] 
& = \E_{\hte_{RW}, Y \mid R, W} 
||\tX \hte_{RW} - \tY||^2 \\ 
& = \E ||\tX \hte_{RW} - \tX \hte_{*} 
+ \tX \hte_{*} - \tY||^2 \\
& = \E ||\tX \hte_{RW} - \tX \hte_{*} 
||^2 + \E || \tX \hte_{*} - \tY||^2 
+ 2 \E (\tX \hte_{RW} - \tX \hte_{*} )^T 
(\tX \hte_{*} - \tY). 
\end{split}    
\end{equation}
The last term is non-positive due 
to the following arguments: 
\begin{equation}
\begin{split}
\E_{\hte_{RW}, Y \mid R, W} 
(\tX \hte_{RW} - \tX \hte_{*} )^T 
(\tX \hte_{*} - \tY)    
& = \E ( \hte_{RW}^T \tX^T \tX \hte_* 
- \hte_{RW}^T \tX^T \tY
- \hte_{*}^T \tX^T \tX \hte_{*} 
+ \hte_{*}^T \tX^T \tY)\\
& =  \E[\hte_{RW}]^T \tX^T \tX \hte_* 
- \E[ \hte_{RW}^T \tX^T \tY]
- \hte_{*}^T \tX^T \tX \hte_{*} 
+ \hte_{*}^T \tX^T \E[\tY]\\
& = - \E[ \hte_{RW}^T \tX^T \tY] 
+ \hte_{*}^T \tX^T \E[\tY]\\
& = - \E[ \tY^T \tX (\tX^T \tX) 
\tX^T \tY] + \E[\tY]^T \tX (\tX^T 
\tX) \tX^T \E[\tY]\\ 
& = - \E[ E^T W^T \tX (\tX^T \tX) 
\tX^T W E]\\
& \leq 0.
\end{split}
\end{equation}
In above, the third line 
applies $\E_{\hte_{RW} \mid R, W}[\hte_{RW}] = \hte_*$ (Claim 
\ref{claim:02}) and the 
fifth line applies 
$\tY = W X \be_* + WE$; 
the last inequality holds 
since $W^T \tX (\tX^T \tX) 
\tX^T W$ is semi-positive 
definite. Plugging 
above back to (\ref{eq:claim03:prf_01})
proves the claim. 
\end{proof}

\begin{claim}
\label{claim:04}
$\E_{\hte_{RW} \mid R,W} 
\frac{1}{n_1} 
|| WXR \hte_{RW} - WXR \hte_* ||^2 
\leq \sigma^2 \frac{p_1}{n_1} 
\lambda_{w \uparrow}^2$.
\end{claim}
\begin{proof}
Let $\tX = WXR$ and $\tY = WY$. 
By design, there is 
\begin{equation}
\tY - \E_{Y \mid R, W}[\tY] 
= (W X \be_* + W E)  
- (W X \be_* + W \E[E]) 
= WE.
\end{equation}
Combining above with 
(\ref{eq:claim02_03}, 
\ref{claim2:02}), we have 
\begin{equation}
\hte_{RW} - \hte_{*} 
= (\tX^T \tX)^{-1} \tX^T 
(\tY - \E_{Y \mid R, W} [\tY]) 
= (\tX^T \tX)^{-1} \tX^T W E.
\end{equation}
Note the left side of this claim 
is $\frac{1}{n_1} 
\E_{\hte_{RW} \mid R,W} 
|| \tX \hte_{RW} - \tX \hte_* ||^2$.  
Plugging in the above, we have 
\begin{equation}
\begin{split}
\frac{1}{n_1} \E_{\hte_{RW} \mid R,W}  
|| \tX \hte_{RW} - \tX \hte_* ||^2 
& = \frac{1}{n_1} \E\ ||\tX (\tX^T \tX)^{-1} \tX^T W E||^2\\ 
& = \frac{1}{n_1} \E\ \mtr ( 
E^T W^T \tX (\tX^T \tX)^{-1} \tX^T
\ \tX (\tX^T \tX)^{-1} \tX^T W E )\\
& = \frac{1}{n_1} 
\E\ \mtr ( \tX (\tX^T \tX)^{-1} \tX^T
\tX (\tX^T \tX)^{-1} \tX^T\ W E E^T W^T )\\
& = \frac{1}{n_1} 
\mtr ( \tX (\tX^T \tX)^{-1} \tX^T
\tX (\tX^T \tX)^{-1} \tX^T W\ 
\E [E E^T]\ W^T )\\
& = \frac{\sigma^2}{n_1} 
\mtr ( \tX (\tX^T \tX)^{-1} \tX^T
\tX (\tX^T \tX)^{-1} \tX^T W W^T )\\
& \leq \frac{\sigma^2}{n_1} 
\lambda^2_{w \uparrow}\ 
\mtr ( \tX (\tX^T \tX)^{-1} \tX^T
\tX (\tX^T \tX)^{-1} \tX^T )\\
& = \frac{\sigma^2}{n_1} 
\lambda^2_{w \uparrow}\ 
\mtr ( I_{p_1} ) = \sigma^2 \frac{p_1}{n_1} 
\lambda^2_{w \uparrow}.
\end{split}
\end{equation}
In above, the fourth line is by 
the linearity of expectation,  
the fifth line is by the fact 
that $\E_{Y \mid R, W} [EE^T] 
= \sigma^2 I_{n}$, and the 
sixth line is by an elementary 
trace inequality $\mtr(B A) 
\leq \lambda_{A \uparrow} 
\mtr(B)$. 
\end{proof}


\begin{claim}
\label{claim:05}
There exists a $W_0 \in \R^{n_1 \times 
n}$ such that
$\E_{R, W, \hte_{RW}} 
[L_{RW}(\hte_{RW})] \geq \frac{n}{n_1} 
\lambda^2_{w_0 \downarrow} 
\E_{R, \hte_{RW}} [L_{R}(\hte_{RW})]$. 
\end{claim}
\begin{proof}
First observe that 
\begin{equation}
\E_{R, W, \hte_{RW}} 
[L_{RW}(\hte_{RW})] 
= \frac{1}{n_1} 
\E_{R, W, \hte_{RW}, Y} 
||WXR \hte_{RW} - WY ||^2 
\geq \frac{1}{n_1} 
\E \lambda^2_{w \downarrow} 
||XR \hte_{RW} - Y||^2
\end{equation}    
By the weighted mean value theorem for integrals, there exists a $W_0 \in \R^{n_1 \times n}$ such that 
\begin{equation}
\frac{1}{n_1} 
\E_{R, W, \hte_{RW}, Y}[   
\lambda^2_{w \downarrow} 
||XR \hte_{RW} - Y||^2] 
= \frac{1}{n_1} 
\lambda^2_{w_0 \downarrow} 
\E_{R, \hte_{RW}, Y} 
||XR \hte_{RW} - Y||^2
= \frac{n}{n_1} 
\lambda^2_{w_0 \downarrow} 
\E_{R, \hte_{RW}} [L_R(\hte_{RW})], 
\end{equation}    
where the last equality 
holds because $L_R(\hte_{RW}) 
= \E_{Y \mid R, \hte_{RW}} 
\frac{1}{n} ||XR \hte_{RW} - Y||^2$.
\end{proof}
 
\begin{remark}
\label{rem:priorkaban}
$\E_{R} [L_{R}(R^T\be_*)] 
\leq L(\be_*) + \frac{1}{p_1} 
||\be_*||_M$, 
where $M = \Sigma + (1+\kappa)\, 
tr(\Sigma) I_p$ and 
$\Sigma = \frac{1}{n} X^T X$.     
\end{remark}
\begin{proof}
This is an intermediate result 
in the proof of \citep[Theorem 1]{kaban2014new}: it was shown 
through (15-17) that 
$L_R(R^T \be_*) 
= L(\be_*) + 
\frac{1}{n} ||X \be_* 
- X R R^T \be_*||^2$ 
and through (19-24) that 
$\frac{1}{n} \E_{R} \left[ ||X \be_* 
- X R R^T \be_*||^2 \right]
\leq \frac{1}{p_1} ||\be_*||
^2_{M}$. Combining both yields 
the remark. (In the original 
proof, $w$ is the true model 
which is $\be_*$ here; $R$ 
is transposed.) 
\end{proof}



\section{Numerical Results}

We empirically tested both 
estimators on real-world data 
and reported their expected 
excess losses and loss gap. 
Since $Y$ is fixed on real-world 
data, all expectations were estimated 
based on random samples of $R$ and $W$, 
and $\be_*$ was estimated as the 
minimizer of $\frac{1}{n}||X \be - Y||^2$. 
All reported results are averaged 
over at least 20 samples of $(R,W)$. 

Here, we present and discuss results 
on the Digit data set; more results 
are in the appendix. The 
Digit set is available on the Python library 
`sklearn.datasets'. It has 1797 
images, each described by 8-by-8 
pixels and containing a digit in 
\{0, \ldots, 9\} (which is the 
domain of $Y$). 
For numerical stability, three pixels 
shared by all images were removed. 

Results are shown in Figure 
\ref{fig:digit}. The left 
subfigure shows expected 
excess loss with $n_1/n = 0.5$, 
where the shadowed 
region represents standard deviation. 
As can be seen, both losses decrease 
as $p_1$ increases, and $\hte_{RW}$ 
generally performs worse than $\hte_{R}$. 
The middle subfigure shows the loss gap 
decreases as $p_1$ increases, and 
seems to slightly bounces back around
$p_1 = 35$. The right subfigure also 
shows the loss gap but based on 
$n_1/n = 0.25$. It can be seen the gap 
bounces back earlier around $p_1 = 25$. 
These observations match the 
theoretical implications discussed 
in Section \ref{sec:mainresult}. 

More results are presented and 
discussed in the appendix. 

\begin{figure}
\centering
\includegraphics[width=.32\linewidth]{figure/Digit_loss.png}
\includegraphics[width=.32\linewidth]{figure/Digit_gap.png}
\includegraphics[width=.32\linewidth]{figure/Digit_gap_smallR.png}
\vspace{-5pt}
\caption{Results on the Digit 
Data Set.}
\label{fig:digit}
\end{figure}




\bibliographystyle{cas-model2-names}
\bibliography{reference}

\newpage 

\newpage 
\appendix 

\section{Additional 
Numerical Results}

We first tested on the Digit data set 
with more choices of $n_1/n$. Figure \ref{fig:digit2} 
shows the results (averaged over 50 random choices 
of $(R,W)$) for $n_1/n \in \{0.05, 0.1, 0.15, 0.3, 0.5, 0.75\}$. 
It is clear that, when $n_1/n$ is smaller, the gap behavior 
transits from decreasing to increasing earlier. When 
$n_1/n$ is as small as $0.05$, the gap only increases. 

We then evaluated both estimators 
on two other real-world data sets 
MNIST\footnote{\url{https://docs.pytorch.org/vision/main/datasets.html}} and 
Superconductivty\footnote{\url{https://doi.org/10.24432/C53P47}}. 
Results are shown in Figure 
\ref{fig:SUP} and Figure 
\ref{fig:MNIST} respectively. 
Similar patterns are 
observed on both sets i.e., 
$\hte_{RW}$ is generally 
worse than $\hte_R$, and 
their gap (as $p_1$ increases) 
first decreases and then 
increases. Then $n_1/n$ 
is smaller, the turnover 
point of their gap is 
also smaller. These observations 
also align with the implications 
of our theoretical results. 





\begin{figure}
\centering
\includegraphics[width=.32\linewidth]{figure/Digit_gap05.png}
\includegraphics[width=.32\linewidth]{figure/Digit_gap10.png}
\includegraphics[width=.32\linewidth]{figure/Digit_gap15.png}
\includegraphics[width=.32\linewidth]{figure/Digit_gap30.png}
\includegraphics[width=.32\linewidth]{figure/Digit_gap50.png}
\includegraphics[width=.32\linewidth]{figure/Digit_gap75.png}
\vspace{-5pt}
\caption{Loss Gap versus $n_1/n$ on the Digit Data Set.}
\label{fig:digit2}
\end{figure}

\begin{figure}
\centering
\includegraphics[width=.325\linewidth]{figure/Figure_super_0.5_1.png}
\includegraphics[width=.325\linewidth]{figure/Figure_super_0.5_2.png}
\includegraphics[width=.325\linewidth]{figure/Figure_super_25.png}
\vspace{-10pt}
\caption{Results on the Superconductivty  
Data Set.}
\label{fig:SUP}
\end{figure}
\begin{figure}
\centering
\includegraphics[width=.325\linewidth]{figure/Figure_mn_0.5_1.png}
\includegraphics[width=.325\linewidth]{figure/Figure_mn_0.5_2.png}
\includegraphics[width=.325\linewidth]{figure/Figure_mn_0.25.png}
\vspace{-10pt}
\caption{Results on the MNIST  
Data Set.}
\label{fig:MNIST}
\end{figure}


   
\end{document}
